\section{Spinor indices and exclusion of unwanted strata}\label{sec:indices}

The purpose of this section is to determine which strata can meet the
cut-down spinor problem.  We separate this numerical assertion from the
regularization needed to apply it: Definition~\ref{def:indices:analytic-hypotheses}
states the precise analytic properties used here.
In particular, none of the dimension arguments below presupposes compactness
of the moduli spaces it is intended to exclude.

\subsection{Conventions and closed indices}

Fix the integral determinant lift $c=c_1(E)$ and its connection, and use only
determinant-one gauge transformations with the prescribed transition maps.
The spin-$c$ determinant is $l=\Lambda-c$.  Write
\begin{equation}\label{eq:indices:charge}
 \kappa=c_2(E)-\frac{c^2}{4},\qquad
 \Theta=\frac{\Lambda^2-\sigma}{4}.
\end{equation}
Here and below a square on a piece with rational homology sphere boundary is
the relative rational square.  On positive spinors the Clifford convention is
$\rho(e^i\wedge e^j)=\rho(e^i)\rho(e^j)$ for orthogonal unit covectors;
positive chirality sees self-dual curvature.  The equations are
\begin{equation}\label{eq:indices:equations}
 D_a\Phi=P_D,\qquad
 \rho(F_a^{0,+})=(\Phi\Phi^*)_{00}+P_+.
\end{equation}
The subscript $00$ means trace-free in both rank-two factors.  The unperturbed
quadratic term has positive pairing with the curvature term in the Dirac
Weitzenb\"ock formula, bounded below by a positive constant times $|\Phi|^4$.
Indeed, represent a spinor by a $2\times2$ complex matrix, diagonalize it by
unitary changes of basis, and project its rank-one Hermitian square off the
two scalar factors.  That projection is nonzero for every nonzero spinor;
its squared norm has a positive minimum on the unit sphere.  This is the
quartic term used in Section~\ref{sec:estimates}.

\begin{definition}\label{def:indices:types}
The phase circle $\mathbb T$ acts by multiplying $\Phi$.  On the quotient by
determinant-one gauge, its kernel at a free orbit is $\{1,-1\}$.  A main
component has type $A$ when $\Phi=0$ and the connection is irreducible, type
$R$ when $\Phi=0$ and the connection has continuous stabilizer, type $S$ when
$\Phi\ne0$ lies in one summand of a parallel splitting, and type $P$
otherwise.  A \emph{fixed-type tuple} has no component of type $P$.
The letters here denote field types; in particular $A$ is not the capped
Seifert surface, and $P$ is not the positive bridge cell.

For a split connection write $E=L_1\oplus L_2$ and
$v=c_1(L_1)-c_1(L_2)$, placing the active spinor line first in type $S$.
Then
\begin{equation}\label{eq:indices:split-classes}
 v\equiv c\pmod2,\qquad K=l+2c_1(L_1)=\Lambda+v.
\end{equation}
At type $R$ choose either ordering, including an ordering at a central
connection.  Actual separated lens caps are not counted as main components
when assigning these types.
\end{definition}

The description of the phase-fixed locus follows directly from stabilizers.
A gauge transformation compensating a phase outside $\{1,-1\}$ is a
parallel unitary automorphism of $E$ with distinct eigenvalues.  If the
spinor is nonzero, it belongs to its corresponding parallel eigenline;
determinant one gives the complementary eigenline.  Conversely,
a parallel splitting carrying the spinor in one summand compensates every
phase by the diagonal determinant-one gauge action.  Thus the above list
also describes the fixed locus of the diagonal phase on a product of main
components.

\begin{lemma}\label{lem:indices:closed}
On a closed filled piece with $b_1=0$, the real instanton index, complex
coupled Dirac index, and real spinor index before phase quotient are
\begin{equation}\label{eq:indices:closed-indices}
 D_I=8\kappa-3(1+b^+),\qquad
 n_D=\Theta-\kappa,\qquad D_{\rm sp}=D_I+2n_D.
\end{equation}
The pure-line tangential index at type $S$ is
\begin{equation}\label{eq:indices:tangent-index}
 d_s=\frac{K^2-(2\chi+3\sigma)}4.
\end{equation}
\end{lemma}
\begin{proof}
The ASD index is the index of $d_a^*\oplus d_a^+$ with
$p_1(\ad E)=-4\kappa$.  For the Dirac operator the degree-four part of
$\operatorname{ch}(E)e^{l/2}\widehat A$ is
\[
 \frac{c^2}{2}-c_2(E)+\frac{cl}{2}+\frac{l^2}{4}
       -\frac{\sigma}{4}
 =\frac{(l+c)^2-\sigma}{4}-\kappa.
\]
Here the signature theorem gives the integrated $\widehat A$ term.
This proves the first two formulas by the index theorem
\cite{AtiyahSingerIII}; the complex Dirac summand contributes twice its
index to the real deformation index.  For a line spinor with determinant
$K$, its complex Dirac index is $(K^2-\sigma)/8$.  Adding twice this to the
scalar gauge/self-duality index $-(1-b_1+b^+)$ gives
\eqref{eq:indices:tangent-index}, since
$2\chi+3\sigma=4(1-b_1+b^+)+\sigma$.
These conventions also agree with
\cite[Equations~(2.51) and~(2.63)]{FeehanLenessLinksI}.
\end{proof}

On cylindrical pieces we use fixed-limit Fredholm operators with a small
positive exponential weight.  The gauge group permits independent parallel
limiting constants in the stabilizers; this is the \emph{unframed}
convention.  Framing an end adds $h^0$, the dimension of its flat
stabilizer.  The true cuts are $J=S^3$ and the lens boundary of a
$(-4)$ disk bundle.  For each flat on these links, $H^1$ of the adjoint local
system is zero: this is immediate on $S^3$, and on its finite lens quotient
it follows by averaging, or by taking invariants in the cohomology of the
finite cover.  The flat Dirac has no kernel by positive scalar curvature.
Thus the weights can be chosen below a common nonzero spectral gap.  Filling
a $J$-end by a flat ball changes no unframed index after stabilizer matching;
the framed ball index is zero.  Across $k$ such cuts the filled ASD indices
sum to $D_I^{\rm tot}-3k$, whereas Dirac indices add without a correction.
These statements follow from the APS index formula with the specified
framing convention \cite{AtiyahPatodiSingerI}.

\subsection{Lens fillings and their indices}

Let $N$ be the disk bundle with zero section $S$, $S^2=-4$.  Normalize the
boundary holonomy by a square root of the determinant on $N$.  The three
states are the central states $+$ and $-$, and the trace-zero state.
This normalization is local to $N$: after a determinant bit is flipped,
the two normalizations need not agree when viewed from the exterior.
An effective cap charge $\kappa_N$ includes particles in the cap and all
trajectory charge in this end, with the boundary state specified as seen
from the main component.

\begin{proposition}\label{prop:indices:lens-index}
The framed ASD and spinor corrections supplied by this end are
\begin{equation}\label{eq:indices:lens-table}
\begin{array}{c|c|c|c}
\text{state}&\text{charge congruence}&\delta_I&\delta_{\rm sp}\\ \hline
 +,-&\kappa_N\in\Z&8\kappa_N&6\kappa_N\\
 \text{trace}&\kappa_N\in\frac14+\Z
 &2+8(\kappa_N-\frac14)&2+6(\kappa_N-\frac14)
\end{array}
\end{equation}
For a split main field there are reference fillings with
\begin{equation}\label{eq:indices:reference-fillings}
\begin{array}{c|c|c|c}
\text{state}&d=v(S)&\kappa_{\rm ref}&\text{choice for an active line}\\ \hline
 +&0&0&K(S)=\Lambda S\\
 -&\pm4&1&|\Lambda S+d|\le4\\
 \text{trace}&\pm2&\frac14&\text{the prescribed boundary eigenline}
\end{array}
\end{equation}
Here $|\Lambda S|=2$.  With these choices the pure-line tangential correction,
including stabilizer matching, is zero.
\end{proposition}
\begin{proof}
A cohomology class evaluating as $r$ on $S$ has square $-r^2/4$.
Consequently a splitting with difference $d$ has charge $d^2/16$.
The normalized eigenline has degree $d/2$ on $S$ and boundary character
$\exp(\pi i d/4)$, up to inversion of the chosen generator.  Thus
$d=0,\pm4,\pm2$ give the three states in
\eqref{eq:indices:reference-fillings}.  Comparing any filling with a
reference by gluing changes the charge by an integer, proving the
congruences.

We compute the local Dirac correction explicitly.  In this calculation
$k(S)$ is the \emph{spin-$c$ determinant evaluation}, not the evaluation of
a twisting line.  On the Hirzebruch surface $\mathbb F_4$, let $T$ be a
fiber divisor.  The canonical class is $-2S-6T$, so the canonical spin-$c$
determinant evaluates as $-2$ on $S$.  Twisting by $\mathcal O(aT)$ gives
$k(S)=-2+2a$.  Projection to $\CP^1$ gives
\[
 \chi(\mathbb F_4,\mathcal O(aT))=a+1.
\]
Contracting $S$ gives $\mathbb P(1,1,4)$, with $\mathcal O(aT)$ corresponding
to $\mathcal O(a)$ on the common smooth complement.  Set
\[
 H(j)=
 \begin{cases}
 \displaystyle\sum_{b=0}^{\lfloor j/4\rfloor}(j-4b+1),&j\ge0,\\
 0,&j<0.
 \end{cases}
\]
Weighted monomials compute $H^0$, intermediate cohomology vanishes, and
Serre duality uses the canonical sheaf $\mathcal O(-6)$.  Hence
\[
 \chi(\mathbb P(1,1,4),\mathcal O(a))=H(a)+H(-a-6);
\]
see \cite[Sections~1.2.3 and~1.4.1]{Dolgachev1982}.
Regard $X=\mathbb P(1,1,4)$ as the complex orbifold with its standard
cyclic quotient chart at the contracted point.  Its orbifold line
$\mathcal L_a=\mathcal O(a)$ has invariant holomorphic-section sheaf
equal to the coherent sheaf denoted $\mathcal O(a)$ above.  On each
finite uniformizing chart, averaging the Dolbeault resolution over
the finite group preserves exactness; invariant partitions of unity
make it a fine resolution.  Thus the orbifold Dolbeault index of
$\mathcal L_a$ equals the coherent Euler characteristic just computed.
The canonical spin-$c$ Dirac operator twisted by $\mathcal L_a$ has
the same symbol and hence the same index
\cite[pp.~155--156]{Kawasaki1979}.

The contraction identifies the complements of $S\subset\mathbb F_4$
and the quotient point of $X$ holomorphically, identifying
$\mathcal O(aT)$ with $\mathcal L_a$.  Since
$K_{\mathbb F_4}=-2S-6T$, this identification also identifies the
canonical spin-$c$ structures on the complement and their boundary
restrictions.  We use this identification, including its boundary
character, for the twisted spin-$c$ structures.  In the quotient
chart, write $\zeta(u,v)=(\zeta u,\zeta v)$ for $\zeta\in\mu_4$,
and choose the character convention
$\mathcal L_a\leftrightarrow\zeta^a$.  The canonical determinant
character is $\zeta^2$, so twisting gives $\zeta^{2+2a}$, consistent
with $k(S)=-2+2a$ modulo four.  The common-complement identification
fixes the full boundary spin-$c$ structure, not merely this
determinant residue.

Choose identical product operators on a separating smooth lens-space
collar.  Cutting with complementary boundary spectral projections
gives the same exterior contribution in the two index problems.
This cutting identity is local on the smooth collar: its parametrix
and matching maps commute with finite-group averaging on the
quotient-ball side.  Consequently only the indices of the two
fillings differ; no characteristic-number formula for an orbifold
with boundary is being assumed.

The quotient-ball filling is $B^4/\mu_4$, with the corresponding
equivariant flat line on $B^4$.  Choose an invariant positive-scalar
cylindrical completion and flat twisting.  The Weitzenb\"ock identity
kills the $L^2$ kernels in both chiralities upstairs, hence in each
required finite-group invariant subspace.  The boundary Dirac
operator is invertible, so a sufficiently small decay weight crosses
no root.  The ball Dirac index is therefore zero.  The common-exterior
comparison gives
\[
 \ind_{\C}D_N^+
 =\chi(\mathbb F_4,\mathcal O(aT))-\chi(X,\mathcal L_a),
\]
which yields the following table:
\begin{equation}\label{eq:indices:dirac-table}
\begin{array}{r|r|r|r|r}
 k(S)&a&\chi(\mathbb F_4,\mathcal O(aT))
 &\chi(\mathbb P(1,1,4),\mathcal O(a))&\ind_{\C}D_N^+\\ \hline
 -6&-2&-1&0&-1\\
 -4&-1&0&0&0\\
 -2&0&1&1&0\\
 0&1&2&2&0\\
 2&2&3&3&0\\
 4&3&4&4&0\\
 6&4&5&6&-1
\end{array}
\end{equation}
For the ASD operator, the framed real tangential line index on $N$ is zero,
since $H^1(N)=0$ and $b^+(N)=0$.  The complex off-diagonal index is obtained
by comparing
\[
 -\chi(\mathbb F_4,\mathcal O(dT))
 -\chi(\mathbb F_4,\mathcal O(-dT))
\]
with the corresponding expression on the weighted plane.  Indeed,
the K\"ahler symbol is the direct sum of the reversed Dolbeault symbols
for the difference line and its inverse, so its complex index is
the negative sum of these two Euler characteristics.  The same
description holds for the corresponding orbifold lines, and the
common-collar comparison applies to these symbols as well.
The flat quotient ball has no deformation cohomology in degrees
one or two.  Its sole unframed contribution is the negative of the
space of parallel gauge generators.  Adding the boundary-framing
factor cancels this contribution, including when the boundary
character is trivial.  Thus its framed index is zero.  The differences
are $0,1,4$ complex dimensions for $d=0,2,4$ respectively.  For example, at
$d=2$ the two resolution Euler characteristics are $3,-1$, and the
orbifold values are $3,0$, giving difference $+1$.  Doubling yields framed
real ASD indices $0,2,8$.

For $\Lambda S=2$, the two Dirac determinant evaluations are $2+d$ and
$2-d$.  Their index sums from \eqref{eq:indices:dirac-table} are $0,0,-1$
at $d=0,2,4$.  Changing the sign of $\Lambda S$ gives the same sums.
Increasing relative charge by one with the boundary data fixed increases
the ASD index by eight and decreases the complex Dirac index by one.
Thus the coupled Dirac corrections are $-\kappa_N$ at a central state and
$\frac14-\kappa_N$ at trace.  Adding twice these to the ASD corrections
proves \eqref{eq:indices:lens-table}.

Finally choose the active-line reference at $-$ with sign opposite to
$\Lambda S$.  The two choices $d=\pm4$ restrict to the same allowable
boundary eigenline: their individual line degrees differ by four.
More explicitly, with the determinant root fixed, the active extensions
differ by a line $M$ of degree four up to sign on $N$, and the complementary
extensions differ by $M^{-1}$.  The restriction
$H^2(N;\Z)\longrightarrow H^2(\partial N;\Z)$ is reduction modulo four.
The reference connection on $M$ therefore has trivial flat boundary
restriction; choose a parallel trivialization there and its inverse on
the complement.  This gives a determinant-one identification preserving
the active boundary line and the fixed transition, and identifies the
spinors themselves, not merely their determinant classes.
Boundary trivializations introduce no $\U(1)$ degree because
$H^1(\partial N;\Z)=0$.
At $+$ and trace the prescribed choices
already give $|K(S)|\le4$.  Equation~\eqref{eq:indices:dirac-table} gives
zero active-line Dirac index in all these cases.  Together with the zero
framed scalar ASD/gauge index, and matching of its parallel $\U(1)$, this
gives zero tangential correction.  Choosing the other reference at $-$
would instead give $|K(S)|=6$ and real tangential correction $-2$; that is
why the choice in \eqref{eq:indices:reference-fillings} is necessary.

The calculation was made with round links.  The positive-scalar link
homotopies in Theorem~\ref{thm:geometry:family} preserve the boundary Dirac
gap; flat adjoint cohomology is unchanged.  Index invariance and excision
with any intervening cylinders therefore give the same answers for the
chosen ends.
\end{proof}

\begin{lemma}\label{lem:indices:lens-charge}
On an actual separated $N$ and its infinite lens cylinders the spinor is
zero.  If its sphere cut is retained, its effective charge satisfies
\begin{equation}\label{eq:indices:lens-charge}
 \kappa_N\in\Z_{\ge1}\quad\text{at }+,-,
 \qquad
 \kappa_N\in\tfrac14+\Z_{\ge0}\quad\text{at trace}.
\end{equation}
At trace charge $1/4$, there are no particles or nonconstant cylinder
levels, and the cap itself has charge $1/4$.
\end{lemma}
\begin{proof}
On $N$ the fixed-line curvature is harmonic ASD and the metric has positive
scalar curvature; the same applies to its exact limiting cylinders with
flat background.  Integrating the Dirac identity with finite-energy
cutoffs therefore annihilates the spinor.  The sufficiently small compact
ASD perturbations allowed on $N$ preserve this conclusion by absorption
in the positive scalar term, with no Dirac perturbation there.
On an unperturbed cylindrical level a nonconstant ASD trajectory has
strictly positive charge, whereas zero charge means a constant flat.
On a cap with a compact curvature perturbation and zero spinor,
Chern--Weil gives
\[
 \kappa_{\rm cap}\ge-C\|F_a^{0,+}\|_{L^2}^2
                 \ge-C'\|P_+\|_{L^2}^2>-\tfrac14
\]
after decreasing the permitted perturbation size.  Its discrete charge
congruence therefore makes its charge nonnegative, and at least $1/4$
when its own end is trace.  Particles carry positive integer charges.
The congruences in
Proposition~\ref{prop:indices:lens-index} now imply the claimed ranges,
except possibly central charge zero.

At $+$, the zero-charge unperturbed cap is flat and its sphere sweep
misses $-1$.  At $-$, a flat extension is impossible because $N$ is simply
connected.  This excludes charge zero with the cut.  These exclusions
persist for the permitted small compact cap perturbations: a sequence
violating them as the perturbation tends to zero has no energy available
for bubbling or nonconstant levels, so converges to the flat case; the
modified sphere path has a uniform avoidance gap there.  The local
compactness and no-loss assertion used in this last sentence is proved
independently of exclusion in Lemmas~\ref{lem:analysis:compactness}
and~\ref{lem:analysis:neck}.

For effective trace charge $1/4$ an integer particle is impossible.  A
nonflat cap with the cut already has charge at least $1/4$ if its own end
is trace, and at least one if it is central.  A zero-charge flat cap at
$+$ misses the cut.  There is consequently no charge left for a
nonconstant cylinder level, proving the final assertion.
\end{proof}

\begin{corollary}\label{cor:indices:trace-minimum}
At the minimal trace reducer, the framed tangential ASD index is zero,
the normal ASD index is one complex dimension, and the coupled Dirac
operator is invertible.  Its normalized eigenline sphere sweep has winding
of absolute value one.
\end{corollary}
\begin{proof}
Order the summands with $d=2$.  The normalized eigenline has degree one.
The harmonic Abelian ASD extension is unique for this ordered boundary
line, since $H^1(N)=0$ and $b^+(N)=0$; any parallel boundary identification
is absorbed by the cap's parallel stabilizer.  The ASD indices were
computed in Proposition~\ref{prop:indices:lens-index}.  The two Dirac
determinant evaluations are $0$ and $\pm4$, so their total index is zero.
Positive scalar curvature and ASD curvature annihilate the positive
Dirac kernel; index zero then annihilates the cokernel.  The degree-one
eigenline gives the asserted winding.  Normal ASD surjectivity and
regularity of the cap cut require the additional choices made in
Section~\ref{sec:gluing}; they are not consequences of this index statement.
\end{proof}

\subsection{Virtual pieces and the analytic hypotheses}

The calculation below applies to a finite stack with nonzero ordinary
count and $n_D^{\rm tot}>0$.  Lemma~\ref{lem:indices:end-exclusion}
specifies the required outside data; with those data fixed,
Theorem~\ref{thm:stack:nonzero} supplies the repetitions.
The ordinary dimension condition and the number of phase sections are
\begin{equation}\label{eq:indices:phase-count}
 D_I^{\rm tot}+n-2n-z=0,\qquad
 \eta=n_D^{\rm tot}-1\ge0.
\end{equation}
There are $n$ interval parameters, $n$ real-degree-two sphere tests, cap
insertions of total real degree $z$, and $\eta$ complex sections of phase
weight two.  The free spinor dimension after these cuts and phase quotient
is therefore
\[
 (D_I^{\rm tot}+n-2n-z)+2n_D^{\rm tot}-2\eta-1=1.
\]

Keep $k$ true $J$-cuts in a limiting tuple, and virtually fill all lens
cuts and all $J$-ends.  The resulting $k+1$ main components are ordered
from left to right.  Let $p_i$ count the remaining interval parameters on
component $i$, restoring a parameter at each of its lens ends.  Let $s_i$
count its assigned sphere tests, and $z_i$ its cap degree.  At a $J$-cut the
one sphere event is assigned to the side on which it occurs.  Allocate
particle charges and integer $J$-trajectory charges to pieces, and include
all effective lens charges in the corresponding virtual filled indices.
Set
\begin{equation}\label{eq:indices:q-and-i}
 q_i=D_{I,i}+p_i-2s_i-z_i,\qquad i_i=q_i+2n_{D,i}.
\end{equation}

\begin{lemma}\label{lem:indices:virtual-sums}
These virtual numbers satisfy
\begin{equation}\label{eq:indices:virtual-sums}
 \sum_i(q_i+4)=4,\qquad
 \sum_i i_i-2\eta=2-4k.
\end{equation}
For a component of type $R$ or $S$, the split reference fillings of
Proposition~\ref{prop:indices:lens-index} give
\begin{equation}\label{eq:indices:ell}
 \kappa_i=-\frac{v_i^2}{4}+\ell_i,\qquad \ell_i\in\Z_{\ge0}.
\end{equation}
\end{lemma}
\begin{proof}
The $J$-gluing formulas and allocation of charges give
\[
 \sum_iD_{I,i}=D_I^{\rm tot}-3k,\quad
 \sum_i n_{D,i}=n_D^{\rm tot},\quad
 \sum_i p_i=n-k,\quad \sum_i s_i=n,\quad \sum_i z_i=z.
\]
Thus $\sum_iq_i=-4k$, proving both sums.  The fourth lost dimension per
$J$ is the removed interval parameter, in addition to the three ASD
stabilizer dimensions.  For a split field,
$c_2(L_1\oplus L_2)-c^2/4=-v^2/4$.  Actual particles and allocated
$J$-trajectory charges add nonnegative integers.  At a lens end the actual
charge minus the chosen reference charge is also a nonnegative integer
by Lemma~\ref{lem:indices:lens-charge} and
\eqref{eq:indices:reference-fillings}.  Their sum is $\ell_i$.
\end{proof}

\begin{definition}\label{def:indices:analytic-hypotheses}
Call a choice of equations and representatives \emph{admissible for the
exclusion calculation} if it has the following properties, for all ideal
main-part tuples in the fixed total charge budget, including the enlarged
candidate class described in Section~\ref{subsec:analysis:candidates}.
In this class the completed main fields satisfy their equations and
retained nonphase cuts, together with the phase cuts when imposed.
Matching across true cuts and the actual lens-cap and connector fields are
omitted.  The assigned lens-cap and connector charges obey their permitted
nonnegative ranges and the total charge constraint; the omitted fields
need not be realizable.
\begin{enumerate}[label=(\roman*)]
\item On a fixed-type tuple, the restricted ASD equations with retained
nonphase cuts on type $A$ components, and the pure tangential equations
on type $S$ components, are transverse.  In particular a type $A$
component has $q_i\ge0$, with
strict inequality if it has any lens end, main particle, or allocated
nonconstant $J$-trajectory loss.  A type $S$ component satisfies
\begin{equation}\label{eq:indices:tangential-bound}
 d_s+p_i+4\ell_i\ge0.
\end{equation}
These are componentwise assertions using only parameters retained on the
component.
\item For types $R,S$, an assigned whole sphere test with $v(S_j)=0$
requires a unit of $\ell_i$.  This includes a separated lens test.  During
switching to the $J$ prescription the unit is required if both half
values vanish; at a true $J$ the assigned test requires the unit if its
own half value vanishes.  Different sphere tests so charged on the same
component require distinct units.  Type $S$ is needed in this assertion
only for a tuple with no $P$.
\item For a tuple containing $P$, project away all type $R$ fields,
positions, and tests, retaining their parameters.  The equations and
retained cuts on the other components are independent of the discarded
data.  Before diagonal phase quotient and phase cuts, the necessary
projected problem has dimension at most
\begin{equation}\label{eq:indices:projection-input}
 \sum_{i\notin R}
 \left(i_i-\sum_{\text{lens ends of }i}(\delta_{\rm sp}-1)-2w_i\right)
       +\sum_{i\in R}p_i,
\end{equation}
where $w_i$ is the total main-particle weight and other allocated integer
loss outside lens caps on that component.  The $\eta$ phase sections
subtract $2\eta$, and the phase quotient subtracts one.  Regularity is
required only near the projections of candidate tuples in the stated
overcount, with all $R$ fields still present before projection.  At those
candidate projections, a negative dimension excludes a zero.  No
assertion is made about other zeros of the enlarged equations after the
$R$ fields are discarded.  This permits finite rational multisections
on free strata without changing the integral instanton maps at zero
spinor.
\end{enumerate}
\end{definition}

Here is the finite-dimensional loss calculation underlying these
requirements.  It also specifies which incidence bounds the analytic
construction must preserve.
\begin{lemma}\label{lem:indices:loss-budget}
For exact sphere, cap surface, and cap point representatives as in
Section~\ref{sec:analysis}, a particle of integer weight $a\ge1$ can
recover at most $4a$ dimensions through its position and any cuts dropped
at that position.  A main particle consequently has net penalty at least
four in ASD dimension and two in spinor dimension.  A separated lens has
net penalty $\delta_I-1$ or $\delta_{\rm sp}-1$ in the respective
cut-down main problem.
\end{lemma}
\begin{proof}
An unconstrained particle position contributes four and releases no cut.
A particle hitting the exact support of a sphere sweep has at most two
position dimensions; dropping the entire sphere test releases two more.
The same budget holds for a cap surface cut, whose position varies in a
two-dimensional domain and whose real degree is two.  At a fixed cap point
there are zero position dimensions and the released degree is four.
Two transverse cap surfaces can be hit at one point: the position is
then zero-dimensional and the released degree is four.  General position
excludes triple surface hits and intersections with the fixed cap points.
Sphere supports are mutually disjoint and disjoint from cap tests.
These cases exhaust the possibilities, also for sphere maps with
self-intersections, because their exact supports are parameterized by
their two-dimensional domains.  Thus the bound is four per distinct
particle position, and at most $4a$ per particle weight.

Dropping the integer charge $a$ lowers $D_I$ by $8a$ and $D_{\rm sp}$ by
$6a$, giving the stated net penalties.  Ignoring matching and using the
same conservative recovery allowance for allocated integer trajectory
units only enlarges the necessary problem.  At a lens end, let $\delta$
be its framed index correction.  Relative to the virtual filled problem,
the main index falls by $\delta$, one parameter is removed, and its
sphere cut of degree two is omitted.  The net change is $-\delta-1+2
=-(\delta-1)$.
\end{proof}

By \eqref{eq:indices:lens-table} and \eqref{eq:indices:lens-charge}, a
lens therefore costs at least one in either calculation.  Its ASD penalty
is at least seven at a central state; its spinor penalty is at least five
there.  At trace the spinor penalty is exactly $1+6h$ when
$\kappa_N=\frac14+h$.  Type $A$ regularity therefore implies (i), with
strictness for every loss.  For type $S$, fill by the sharp tangential
references and ignore all retained cuts; allowing four position
dimensions per unit of $\ell_i$ gives
\eqref{eq:indices:tangential-bound}.  The remaining issue is to obtain
these regularity and independence assertions simultaneously.
Section~\ref{sec:analysis} proves them, rather than assuming regularity of
a split-component Dirac operator with negative index.

\subsection{The square estimate on a positive cell}

We now use the lattice of Proposition~\ref{prop:stack:lattice} and the
metric tests of Theorem~\ref{thm:estimates:clamp}.  In a positive cell put
\[
 x=vG,\qquad u=vU,\qquad t_b=vT_b\quad(1\le b\le4).
\]
The cell contributes
\begin{equation}\label{eq:indices:positive-square}
 Q_0=2xu-2u^2-\sum_{b=1}^4t_b^2
\end{equation}
to $v^2$.  Exactly two $t_b$ are odd.  The two neighboring negative half
values will be denoted $h_1,h_2$, and the whole side evaluations are
\begin{equation}\label{eq:indices:d-values}
 d_1=x-\sum_bt_b-h_1,\qquad d_2=x-2u-h_2.
\end{equation}
All $x,h_i,d_i$ are even.  In particular, for a selected nonempty subset
$I$ of whole sides,
\begin{equation}\label{eq:indices:z-local}
 z_I=u+\frac14\sum_{i\in I}d_i\in\tfrac12\Z.
\end{equation}
We retain the subscript here to distinguish $z_I$ from the cap insertion
degree $z$.

For endpoint bits $e,f\in\{0,1\}$, the background values are
\begin{equation}\label{eq:indices:lambda-values}
 \lambda=-1-e+f,\quad
 \lambda_{\{1\}}=-\tfrac32+f,\quad
 \lambda_{\{2\}}=-\tfrac12-e,\quad
 \lambda_{\{1,2\}}=-1.
\end{equation}
If $u\ne0$, the clamp selects a nonempty subset of available \emph{whole}
sides with
\begin{equation}\label{eq:indices:clamp-input}
 \operatorname{sign}(u)z_I\le0\quad(R),\qquad
 \operatorname{sign}(u)z_I<|\lambda_I|\quad(S).
\end{equation}
A side whose test has switched to the separate-half rule is not needed
in this subset.  A sphere in a separated lens cap still counts as whole
for this assertion and for its reference filling.

\begin{lemma}\label{lem:indices:positive-square}
For a cell of type $R$ or $S$, combine its square with the neighboring
negative half contributions $-\frac12\sum_{i\in I}h_i^2$ selected by the
clamp.  The resulting square $Q$ is at most $-2$, except possibly in type
$S$ when $I$ is a singleton and $u=z_I=1$ or $u=z_I=-1$.
In that exception $Q\le2$ and the selected whole side has $d_i=0$.
If $u=0$, the cell square alone is at most $-2$.
\end{lemma}
\begin{proof}
For $u=0$, Equation~\eqref{eq:indices:positive-square} and the two odd
$t_b$ give $Q_0\le-2$.  Suppose $u\ne0$.  Direct substitution from
\eqref{eq:indices:d-values}--\eqref{eq:indices:z-local} gives
\begin{align}
 I=\{1\}:\quad Q
 &=8z_Iu-4u^2-\sum_b(t_b-u)^2-\frac{(h_1-2u)^2}{2},
 \label{eq:indices:square-one}\\
 I=\{2\}:\quad Q
 &=8z_Iu-4u^2-\sum_bt_b^2-\frac{(h_2-2u)^2}{2},
 \label{eq:indices:square-two}\\
 I=\{1,2\}:\quad Q
 &=4z_Iu-2u^2-\sum_b(t_b-u/2)^2
                     -\sum_{i=1}^2\frac{(h_i-u)^2}{2}.
 \label{eq:indices:square-both}
\end{align}
Let $a=|u|\ge1$ and $w=\operatorname{sign}(u)z_I$.
For a singleton the exceptional-square sum is at least two, also after
shifting each $t_b$ by the integer $u$: exactly two of the shifted values
remain odd.  Type $R$ has $w\le0$ and therefore $Q\le-2$.
For type $S$, the singleton threshold is $1/2$ or $3/2$.
Since $w\in\frac12\Z$, the former implies $w\le0$ and the latter
$w\le1$.  Thus
\[
 Q\le8aw-4a^2-2\le8a-4a^2-2.
\]
For $a\ge2$ the last expression is at most $-2$.  For $a=1$, every
$w\le1/2$ still gives $Q\le-2$.  The remaining possibility is $a=1,w=1$,
that is, $u=z_I=\pm1$.  Its upper bound is two, and
$d_i=4(z_I-u)=0$.

For two sides, type $R$ again has $w\le0$, while type $S$ has $w<1$,
hence $w\le1/2$.  If $u$ is even, the $t_b-u/2$ are integers with two
odd values, so their squared sum is at least two.  If $u$ is odd, all
four are half-integers, contributing at least one; both $h_i-u$ are odd,
so the final sum in \eqref{eq:indices:square-both} contributes at least
one more.  In both cases
\[
 Q\le2a-2a^2-2\le-2.
\]
This proves every case, including the stated exceptional alternative.
\end{proof}

The exception requires an actual loss unit.  Its side was selected among
whole tests, so property (ii) of
Definition~\ref{def:indices:analytic-hypotheses} applies to $d_i=0$.
This use of a whole side is the reason the metric construction stops using
a side before its sphere representative switches to the $J$ rule.

\subsection{Incidence assignments and fixed-type runs}

Consider an interior filled component between true $J$-cuts.  Let $L_i$
be its number of cells, and $m_i$ its number of positive cells.  By
Equation~\eqref{eq:stack:interior-topology},
\begin{equation}\label{eq:indices:interior-data}
 p_i=L_i-1,\qquad b_i^+=m_i,\qquad
 \Theta_i=m_i+\frac{\Delta_i}{2}.
\end{equation}
The bit differences $\Delta_i$ telescope on contiguous runs, so their
sum is the difference of two bits and belongs to $\{-1,0,1\}$.
Each negative cell contributes minus one half the sum of the squares of
its two even half-sphere evaluations.

We specify the assignment convention carefully.  Reserve two potential
slots for each positive cell, one at each adjacent test.  A slot is
occupied by the whole test or by the half test assigned to that cell's
side of a $J$-break.  Let $s_P\le2m_i$ be the number
of the component's assigned tests occupying these slots, and put
$s_N=s_i-s_P$.  All other assigned tests use the negative rule: charge one
nonzero even half value in the component, or, if all the relevant half
values vanish, charge a unit of $\ell_i$.  A nonzero even half contributes
at most $-2$ to $v^2$.

The negative halves used with the clamp in
Lemma~\ref{lem:indices:positive-square} are reserved along with the
positive-cell slots, and are never used by the negative rule.  Each half
belongs to one adjacent test, so no half is charged twice.  Positive
cells have disjoint adjacent macros at the prescribed spacing.  A test
assigned across a $J$ to a neighboring negative component uses the
negative half in that component; the positive-cell slot on the other side
is then empty.  This convention also applies when only one of those
components belongs to a run under consideration.  At an outside cap,
the bordering negative trace half is available to the negative rule.
During switching the negative rule only needs a nonzero half or the
incidence unit when both half values vanish, exactly as in
Definition~\ref{def:indices:analytic-hypotheses}.

\begin{lemma}\label{lem:indices:component-score}
For an interior fixed component of type $R$ or $S$ satisfying the
incidence property (ii) of
Definition~\ref{def:indices:analytic-hypotheses},
\begin{equation}\label{eq:indices:charge-lower}
 8\kappa_i\ge\max(4s_i-4m_i,\,4m_i),
\end{equation}
and
\begin{equation}\label{eq:indices:q-score}
 q_i+4\ge
 \max(2s_i+L_i-7m_i,\,m_i+L_i-2s_i).
\end{equation}
More precisely, if $z_0$ negative-rule assignments have no nonzero half
option and $e_0$ positive cells have the exception of
Lemma~\ref{lem:indices:positive-square}, then
\begin{equation}\label{eq:indices:incidence-charge}
 \ell_i\ge z_0+e_0,\qquad
 8\kappa_i\ge
 4(s_N-z_0)+4m_i-8e_0+8\ell_i
 \ge4s_N+4m_i.
\end{equation}
\end{lemma}
\begin{proof}
The zero negative-rule tests and the exceptional whole tests are distinct
assigned tests with disjoint sphere supports.  Property (ii) therefore
charges distinct units and gives $\ell_i\ge z_0+e_0$.
The negative rule contributes at most $-2(s_N-z_0)$ to $v_i^2$.
The positive-cell estimates contribute at most $-2m_i+4e_0$.
Every remaining negative square can be discarded in an upper bound.
Multiplying by $-2$ in \eqref{eq:indices:ell} gives the middle inequality
of \eqref{eq:indices:incidence-charge}.  Substitution of
$\ell_i\ge z_0+e_0$ gives the last one (in fact it leaves an additional
$4z_0$).  Now $s_N\ge s_i-2m_i$ and $s_N\ge0$, proving
\eqref{eq:indices:charge-lower}.  Finally, since there are no cap
insertions on an interior piece,
\[
 q_i+4=8\kappa_i+L_i-3m_i-2s_i.
\]
Using the two charge bounds gives \eqref{eq:indices:q-score}.
\end{proof}

\begin{lemma}\label{lem:indices:fixed-run}
For any nonempty contiguous run of interior components of types $R,S$,
\begin{equation}\label{eq:indices:fixed-run}
 \sum_{i\text{ in the run}}(q_i+4)\ge-2.
\end{equation}
\end{lemma}
\begin{proof}
Let $L,m,s$ be the totals on the run.  Summing each of the two separate
bounds in \eqref{eq:indices:q-score} gives
\[
 \sum(q_i+4)\ge\max(2s+L-7m,\,m+L-2s).
\]
Every internal adjacency contributes one assigned test to the run, while
only its two endpoints can contribute an additional test.  Thus
$L-1\le s\le L+1$.  Corollary~\ref{cor:stack:interior} gives
$L\ge4m-3$ for $m\ge1$.
If $m\ge2$, the first bound is at least
$3L-7m-2\ge5m-11\ge-1$.  For $m=1,L\ge3$ it is at least $3L-9\ge0$.
The remaining short possibilities are
\[
\begin{array}{c|c|c}
 (m,L)&\text{possible }s&
 \max(2s+L-7m,\,m+L-2s)\\ \hline
 (1,1)&0,1,2&2,0,-2\\
 (1,2)&1,2,3&1,-1,1.
\end{array}
\]
For $m=0$ the first bound is at least $L>0$.  Every case is at least $-2$.
\end{proof}

\subsection{The outside caps and the large exceptional evaluations}

The constants in Proposition~\ref{prop:estimates:cap} depend only on the
old caps $W$, before the extra blow-ups.  They are independent of both the
new exceptional evaluations and the number of intervals.  We use both
of its square bounds, one for $v$ and one for $K$.

\begin{lemma}\label{lem:indices:end-exclusion}
The delays before the first and after the last positive bridge, followed
by the odd values $\Lambda(E_{\rm exc})$, can be chosen so that no outside
component has type $R$, and every outside component of type $S$ satisfying
properties (i) and (ii) of
Definition~\ref{def:indices:analytic-hypotheses} has
\begin{equation}\label{eq:indices:end-score}
 q_i+4\ge8.
\end{equation}
This holds also for a single component containing both outside caps.
The choices precede the number of middle repetitions and all finite
metric isolation sizes.
\end{lemma}
\begin{proof}
The exclusion of $R$ is part of Proposition~\ref{prop:estimates:cap}.
Consider type $S$.  Let $L$ count its internal cells and $m$ its positive
cells.  Rational orthogonal splitting separates the old cap parts, the
new exceptional summands, the bordering negative trace halves, and the
internal cells.  Applying the incidence argument of
Lemma~\ref{lem:indices:component-score}, including the bordering trace
halves, gives
\begin{equation}\label{eq:indices:end-square}
 8\kappa_i\ge4s_i-4m
       +2\sum_{E_{\rm exc}\text{ in }i}v(E_{\rm exc})^2-C_1.
\end{equation}
To see that $C_1$ is uniform, the old-cap terms in $v^2$ are bounded above
by the sum of their fixed constants $C_W$, the exceptional summands
contribute $-v(E_{\rm exc})^2$, and all unused negative squares can only
improve the estimate.  The same $z_0,e_0$ incidence argument cancels the
possible positive-cell exceptions.  Only one or two fixed cap pieces
are involved.

On an end component $p_i=L+O(1)$, $s_i\ge L-O(1)$, and
$b_i^+=m+O(1)$; all the constants, including $z_i$, depend only on the
fixed caps.  Substituting \eqref{eq:indices:end-square} in
\eqref{eq:indices:q-and-i} therefore gives a single fixed $C$ with
\begin{equation}\label{eq:indices:end-lower}
 q_i\ge3L-7m
       +2\sum_{E_{\rm exc}\text{ in }i}v(E_{\rm exc})^2-C.
\end{equation}
For example, the intermediate expression is
$2s_i+L-7m+2\sum v(E_{\rm exc})^2-C'$; using $s_i\ge L-O(1)$
gives the displayed form.  Spacing gives $m\le(L+3)/4$, so its first two
terms are at least $(5L-21)/4$.  Choose a fixed $L_*$ so large that
$L\ge L_*$ already implies $q_i\ge4$.

Delay the first and last positive cells beyond $L_*$.  An end component
with $L<L_*$ then has no positive cell.  Suppose, seeking the only
remaining obstruction, that $q_i<4$.  All its topological dimension terms
are bounded independently of the new values of $\Lambda(E_{\rm exc})$,
so \eqref{eq:indices:q-and-i} bounds $\kappa_i$ above.  There are no
positive cells, and the cap estimate for $v$ gives $v_i^2\le C_v$,
independently of those values.  Thus
\[
 0\le\ell_i=\kappa_i+\frac{v_i^2}{4}\le C_\ell.
\]
The tangential necessary condition
\eqref{eq:indices:tangential-bound} now gives
\[
 K_i^2\ge2\chi_i+3\sigma_i-4p_i-16\ell_i\ge-C_K.
\]
There are only boundedly many negative cells and fixed caps in these
short end components, so $C_K$ is uniform.  The \emph{second} cap estimate,
for $K$, and negativity of all the other summands give
\[
 K_i^2\le C'_K-\sum_{E_{\rm exc}\text{ in }i}K(E_{\rm exc})^2.
\]
It follows that every $|K(E_{\rm exc})|$ is bounded by a fixed number $B$,
independent of $\Lambda(E_{\rm exc})$.  But
\[
 |v(E_{\rm exc})|
 =|K(E_{\rm exc})-\Lambda(E_{\rm exc})|
 \ge|\Lambda(E_{\rm exc})|-B.
\]
Choose the odd exceptional evaluations at both ends so large that
\eqref{eq:indices:end-lower} is incompatible with $q_i<4$ even for
$L<L_*$.  A component containing either end contains its exceptional
summand, so the argument covers both single-cap and two-cap components.
This proves \eqref{eq:indices:end-score}.

The selection just made uses only the old-cap constants and the bounded
short-end topologies: first $L_*$ and the delays, then the exceptional
values.  The middle repetitions can subsequently be chosen to make
$n_D^{\rm tot}>0$ and the ordinary count nonzero.  Only after that finite
problem is fixed are the isolation lengths and perturbation sizes chosen
as in Corollary~\ref{cor:estimates:choices}.  In particular the cap bounds
are not being asserted uniformly over unboundedly many new field
problems at one fixed isolation length.
\end{proof}

\begin{proposition}\label{prop:indices:fixed-exclusion}
With the preceding cap choices, a fixed-type tuple satisfying properties
(i) and (ii) of Definition~\ref{def:indices:analytic-hypotheses} has no
true $J$-cut.  No mixed projection regularity is assumed in this assertion.
It is a single type $A$ component with no lens end, particles, or other
positive loss.  Its cut-down instanton dimension is zero.
\end{proposition}
\begin{proof}
Suppose first that $k>0$, so the two outside components are distinct.
Each is type $A$ or $S$.  Let $a$ be the number of all type $A$ components,
and $b\in\{0,1,2\}$ the number of type $S$ outside components.
There are $a+b\ge2$ distinguished components.  They contribute at least
$4a+8b$ to $\sum(q_i+4)$.  The remaining components form at most
$a+b-1$ contiguous interior runs of types $R,S$, each contributing at
least $-2$ by Lemma~\ref{lem:indices:fixed-run}.  Thus
\[
 \sum_i(q_i+4)\ge4a+8b-2(a+b-1)
                  =2a+6b+2\ge6,
\]
contradicting \eqref{eq:indices:virtual-sums}.  A run adjacent to an
outside $S$ uses only its interior portion in this count; the outside
component has already contributed its separate score of eight.

If $k=0$, there is a single component and $q_1=0$.  It cannot be type $R$
by the cap exclusion, or type $S$ by \eqref{eq:indices:end-score}.
It is therefore type $A$.  Property (i) rules out every lens end,
particle, and allocated trajectory loss, since any of these would force
$q_1>0$.  Equation~\eqref{eq:indices:phase-count} is exactly its ordinary
cut-down dimension-zero condition.
\end{proof}

\subsection{Projection away from reducible zero-spinor components}

For the mixed calculation only type $R$ is discarded.  The sharper square
bound for that type has no exceptional positive-cell case.

\begin{lemma}\label{lem:indices:mixed-run}
For a nonempty contiguous run of $r$ interior components of type $R$,
\begin{equation}\label{eq:indices:mixed-run}
 \sum_{i\text{ in the run}}(4+i_i-p_i)\ge-2.
\end{equation}
More precisely, let $s,m$ be the total assigned tests and positive cells,
let $s_P$ count occupied positive-cell slots, and let
$\mathrm{miss}=2m-s_P$.  With $\Delta$ the telescoping endpoint bit
difference,
\begin{equation}\label{eq:indices:mixed-run-sharp}
 \sum(4+i_i-p_i)
 \ge s-4m+r+\Delta+3\,\mathrm{miss}.
\end{equation}
\end{lemma}
\begin{proof}
From the definitions and \eqref{eq:indices:interior-data},
\[
 i_i-p_i=6\kappa_i-3-m_i+\Delta_i-2s_i.
\]
With $e_0=0$, Equation~\eqref{eq:indices:incidence-charge} gives
$6\kappa_i\ge3s_{N,i}+3m_i$.  Consequently
\[
 4+i_i-p_i\ge
 1+s_i-4m_i+\Delta_i+3(2m_i-s_{P,i}),
\]
which proves \eqref{eq:indices:mixed-run-sharp} after summing.

If $m=0$, its right side is $s+r+\Delta\ge0$ because
$r\ge1$ and $\Delta\ge-1$.  Suppose $m>0$.  Write $a,b\ge0$ for the
numbers of cells before the first and after the last positive cell in
the run, and $\epsilon_a,\epsilon_b\in\{0,1\}$ for the two endpoint
test assignments into the run.  Then
\[
 L\ge4m-3+a+b,\qquad s=L-1+\epsilon_a+\epsilon_b.
\]
An endpoint with zero offset and no assignment leaves its boundary
positive-cell slot empty.  It therefore contributes one to
$\mathrm{miss}$.  At the left endpoint,
\[
 a+\epsilon_a+3\,\mathbf1_{\{a=0,\epsilon_a=0\}}\ge1,
\]
and the identical statement holds at the right endpoint.  The two
empty boundary slots are distinct, including when there is only one
positive cell.  Applying these inequalities to
\eqref{eq:indices:mixed-run-sharp} gives
\[
 \sum(4+i_i-p_i)
 \ge r+\Delta-4+a+\epsilon_a+b+\epsilon_b+3\,\mathrm{miss}
 \ge r+\Delta-2\ge-2.
\]
When a boundary test is assigned to a negative component outside the run,
its cost is charged there by its negative half; inside the run the empty
positive slot gives precisely the correction just used.  No charge from
an exterior component is borrowed in this calculation.
\end{proof}

\begin{proposition}\label{prop:indices:mixed-exclusion}
A tuple containing type $P$, satisfying the type $R$ incidence in
property (ii) and the projection property (iii) of
Definition~\ref{def:indices:analytic-hypotheses}, has no true $J$-cut.
It has no particles or other integer losses.  Apart from an unbroken
free main field, the only possible case is one separated lens end of
effective trace charge $1/4$, a single main component of type $P$, and no other
lens end or loss.
\end{proposition}
\begin{proof}
Write $N$ for the number of non-$R$ main components, and
$N_{\mathrm{lens},i}$ for the number of lens ends on component $i$.
Since every lens penalty in
\eqref{eq:indices:projection-input} is at least one, the necessary
projected dimension after phase and phase cuts is at most
\begin{align}
 &\sum_{i\notin R}(i_i-N_{\mathrm{lens},i}-2w_i)
       +\sum_{i\in R}p_i-2\eta-1
 \notag\\
 &\qquad=5-4N-\sum_{i\in R}(4+i_i-p_i)
                    -\sum_{i\notin R}(N_{\mathrm{lens},i}+2w_i).
 \label{eq:indices:projected-dimension}
\end{align}
Indeed substitute $\sum i_i-2\eta=2-4k$ and
$k+1=N+\#R$ to obtain the equality.  This is a dimension of a necessary
projection, so forgetting interpiece matching can only enlarge the
candidate space.  Its negative value excludes the original tuple
without any regularity assertion for the discarded $R$ fields.

If $k>0$, both outside components are non-$R$ by
Lemma~\ref{lem:indices:end-exclusion}, and hence $N\ge2$.
There are at most $N-1$ contiguous $R$-runs.  Applying
Lemma~\ref{lem:indices:mixed-run} in
\eqref{eq:indices:projected-dimension} bounds it above by
\[
 5-4N+2(N-1)=3-2N\le-1.
\]
Thus $k=0$.  There is only one main component, and because the tuple
contains $P$ it has type $P$.  The projected dimension is at most
$1-N_{\mathrm{lens},1}-2w_1$, so $w_1=0$ and
$N_{\mathrm{lens},1}\le1$.  If there is a lens end, use the actual penalty
$\delta_{\rm sp}-1$ in \eqref{eq:indices:projection-input}.  A central
cap costs at least five, while a trace cap of charge $\frac14+j$ costs
$1+6j$.  Nonnegative dimension therefore forces trace charge exactly
$1/4$.  For an actual limiting tuple,
Lemma~\ref{lem:indices:lens-charge} also excludes particles or
breaks inside that cap and its cylindrical end.  These are all the
possibilities asserted.
\end{proof}

\begin{theorem}[Numerical exclusion]\label{thm:indices:exclusion}
Choose the old-cap data, outside delays, exceptional evaluations, finite
stack, metric sizes, and sufficiently small equation data in the order
of Lemma~\ref{lem:indices:end-exclusion} and
Corollary~\ref{cor:estimates:choices}, with $n_D^{\rm tot}>0$.
Assume the regularity, incidence, and projection properties of
Definition~\ref{def:indices:analytic-hypotheses}.  Then the only fixed-type tuples
meeting the nonphase cuts are the legitimate unbroken isolated type $A$
instantons.  The free part after the $\eta=n_D^{\rm tot}-1$ phase cuts
and phase quotient has expected dimension one.  Its only possible
true-face or positive-loss limiting tuples are single trace-minimum
lens faces, with one type $P$ main component and no other losses.
Consequently the only boundary types to be constructed for its
compactified one-dimensional count are links of the legitimate instantons
and those trace-minimum ends.
\end{theorem}
\begin{proof}
Proposition~\ref{prop:indices:fixed-exclusion} gives the fixed-type
assertion, and Equation~\eqref{eq:indices:phase-count} gives the dimension.
Proposition~\ref{prop:indices:mixed-exclusion} gives the asserted free
degenerations.  Lemma~\ref{lem:indices:lens-charge} describes the minimal
cap.  The analytic realization of these exclusions and of the instanton
links is proved in Section~\ref{sec:analysis}; regular trace-end gluing
and orientation comparison are proved in Section~\ref{sec:gluing}.
\end{proof}
